% Modified Template by Jonathan Doucette and Kevin Multani, original by: Jonathan Ward

\documentclass[12pt]{article} 
\usepackage[english]{babel}
\usepackage[utf8]{inputenc}
\usepackage{amsmath} % AMS Math Package
\usepackage{amsthm} % Theorem Formatting
\usepackage{amssymb} % Math symbols such as \mathbb
\usepackage{graphicx} % Allows for eps images
\usepackage{multicol} % Allows for multiple columns
\usepackage[dvips,letterpaper,margin=1in,bottom=1in]{geometry}
\usepackage{hyperref}
\usepackage{parskip} % Removes indentation from paragraphs
\usepackage{xcolor,xspace,soul} % Colour, spacing, and highlighting
\usepackage{mathrsfs}
\usepackage{bm} % For bold math symbols
\usepackage{amscd}
\usepackage[all,cmtip]{xy}
%\usepackage{bbm}
\usepackage{titling}
\usepackage{listing} % for code snippets
%\usepackage{minted} % for code snippets
\usepackage{enumerate}
\usepackage{fancyhdr}
\usepackage[]{physics}
\usepackage[makeroom]{cancel}
\usepackage{pdfpages}
\usepackage[]{mcode}
\usepackage[title]{appendix}

% ***********************************************************
% ********************** BEGIN TITLE PAGE *******************
% ***********************************************************
\newcommand*{\titleGM}{\begingroup % Create the command for including the title page in the document
\hbox{ % Horizontal box
\hspace*{0.2\textwidth} % Whitespace to the left of the title page
\rule{1pt}{\textheight} % Vertical line
\hspace*{0.05\textwidth} % Whitespace between the vertical line and title page text
\parbox[b]{0.75\textwidth}{ % Paragraph box which restricts text to less than the width of the page

{\noindent\Huge\bfseries MATH 521}\\[2\baselineskip] % Title
{\large \textit{Assignment 4}}\\[4\baselineskip]
%{\large \textsc{ Jonathan Doucette }} % Author name

\vspace{0.5\textheight} % Whitespace between the title block and the publisher
{\noindent \today }\\[\baselineskip]
%{\noindent Student Number: 35298124 }\\[\baselineskip] 
%{\noindent \footnotesize All problems below are \textit{Copyright \copyright \space 2018 Timm Treskatis. All Rights Reserved.} }\\[\baselineskip]
} % end parbox
} % end hbox
\endgroup}

 % Sets margins and page size
\pagestyle{fancy} 

%\lhead{Jonathan Doucette}
\rhead{\today}
\rfoot{Page \thepage}
\cfoot{}

\makeatletter % Need for anything that contains an @ command 
\renewcommand{\maketitle} % Redefine maketitle to conserve space
{ \begingroup \vskip 10pt \begin{center} \Huge {\bf \@title}
\vskip 10pt \large \@author \hskip 20pt \@date \end{center}
  \vskip 10pt \endgroup \setcounter{footnote}{0} }
\makeatother % End of region containing @ commands

% ***********************************************************
% ********************** END TITLE PAGE *********************
% ***********************************************************

% ***********************************************************
% ********************** BEGIN NEW COMMANDS *****************
% ***********************************************************

\renewcommand{\labelenumi}{(\alph{enumi})} % Use letters for enumerate
\let\vaccent=\v % rename builtin command \v{} to \vaccent{}

%% MISC
\newcommand{\ab}[1]{\left| #1 \right|} % for absolute value
\newcommand{\avg}[1]{\left< #1 \right>} % for average
\let\underdot=\d % rename builtin command \d{} to \underdot{}
\let\baraccent=\= % rename builtin command \= to \baraccent
\renewcommand{\=}[1]{\stackrel{#1}{=}} % for putting numbers above =
\providecommand{\fr}{\frac}
\providecommand{\RR}{\mathbb{R}}
\providecommand{\CC}{\mathbb{C}}
\providecommand{\NN}{\mathbb{N}}
\providecommand{\e}{\epsilon}
\DeclareMathOperator{\di}{d\!}
\newcommand*\ieval[3]{\left.#1\right\rvert_{#2}^{#3}}

%% Vectors
\renewcommand{\v}[1]{\ensuremath{\mathbf{#1}}} 
\newcommand{\gv}[1]{\ensuremath{\mbox{\boldmath$ #1 $}}} % for vectors of Greek letters
\newcommand{\uv}[1]{\ensuremath{\mathbf{\hat{#1}}}} % for unit vector
\providecommand{\wave}[1]{\v{\tilde{#1}}}

%% DERIVATIVES
\renewcommand{\d}[2]{\frac{d #1}{d #2}} % for derivatives
\newcommand{\dubd}[2]{\frac{d^2 #1}{d #2^2}} % for double derivatives
\newcommand{\pd}[2]{\frac{\partial #1}{\partial #2}} % for partial derivatives
\newcommand{\pdd}[2]{\frac{\partial^2 #1}{\partial #2^2}} % for double partial derivatives

%% Operators
\newcommand{\Gradient}{\ensuremath{\mbox{\boldmath$\nabla$}}} % gradient

%% Text
\newcommand{\mathcolorbox}[2]{\colorbox{#1}{$\displaystyle #2$}}
\newcommand{\hlfancy}[3]{\textcolor{#1}{\sethlcolor{#2}\hl{#3}}}
\newcommand{\TODO}[1]{\hlfancy{red}{yellow}{\textbf{TODO: #1}}}

%% Code
%\newcommand{\code}[1]{\mintinline{C}{#1}}
%\newcommand{\code}[1]{\texttt{#1}}
\newcommand{\code}[1]{\lstinline[columns=fixed]{#1}}
\newcommand{\includecode}[1]{\lstinputlisting{#1}}

% ***********************************************************
% ********************** END NEW COMMANDS *******************
% ***********************************************************

% ***********************************************************
% ********************** BEGIN NEW ENVS *********************
% ***********************************************************

% Theorem
\newenvironment{theorem}[2][Theorem]{\begin{trivlist}
\item[\hskip \labelsep {\bfseries #1}\hskip \labelsep {\bfseries #2.}]}{\end{trivlist}}
% Lemma
\newenvironment{lemma}[2][Lemma]{\begin{trivlist}
\item[\hskip \labelsep {\bfseries #1}\hskip \labelsep {\bfseries #2.}]}{\end{trivlist}}
% Corollary
\newenvironment{corollary}[2][Corollary]{\begin{trivlist}
\item[\hskip \labelsep {\bfseries #1}\hskip \labelsep {\bfseries #2.}]}{\end{trivlist}}

% Exercise
\newenvironment{exercise}[2][Exercise]{\begin{trivlist}
\item[\hskip \labelsep {\bfseries #1}\hskip \labelsep {\bfseries #2.}]}{\end{trivlist}}
% Problem
\newenvironment{problem}[2][Problem]{\begin{trivlist}
\item[\hskip \labelsep {\bfseries #1}\hskip \labelsep {\bfseries #2.}]}{\end{trivlist}}
% Question
\newenvironment{question}[2][Question]{\begin{trivlist}
\item[\hskip \labelsep {\bfseries #1}\hskip \labelsep {\bfseries #2.}]}{\end{trivlist}}
% Solution
\newenvironment{solution}{\begin{proof}[Solution]}{\end{proof}}

% Afterword
\newenvironment{afterword}[2][Appendix]{\begin{trivlist}
\item[\hskip \labelsep {\bfseries #1}\hskip \labelsep {\bfseries #2}]}{\end{trivlist}}

% ***********************************************************
% ********************** END NEW ENVS ***********************
% ***********************************************************

% ***********************************************************
% ********************** END TITLEPAGE **********************
% ***********************************************************
%\documentclass[10pt,letterpaper]{scrartcl}
%\usepackage{amsfonts,amsmath,amssymb,braket,xcolor,dsfont,enumerate,fontawesome,graphicx}
\usepackage{amsfonts,amsmath,amssymb,braket,xcolor,enumerate,graphicx}
\usepackage[hidelinks]{hyperref}
\usepackage{listings,multicol,mathtools,textcomp,tikz,pgfplots,wrapfig}
\usepackage[inner=2cm,outer=2cm,top=2cm,bottom=2cm]{geometry}
\usepackage{tabularx}
\usepackage{booktabs}
\usepackage{appendix}
\usepackage{fancyhdr}
\usetikzlibrary{arrows}
%\pgfplotsset{compat=1.12}

\pagestyle{empty}
\setlength{\parindent}{0pt}
\setlength{\parskip}{6pt}

\newcommand{\dx}{\;\mathrm{d}x}
\newcommand{\ds}{\;\mathrm{d}s}

\begin{document}

\begin{minipage}{.2\textwidth}
\includegraphics[width=42pt]{ubc-logo.png}
\end{minipage}
\hfill
\begin{minipage}{.75\textwidth}
\setlength{\parskip}{6pt}
\begin{flushright}
{\sffamily
\textbf{MATH521}\\
\textbf{Numerical Analysis of Partial Differential Equations}

Winter 2017/18, Term 2\\
Timm Treskatis
}
\end{flushright}
\end{minipage}

\section*{Homework Assignment 12}

%Please submit the following files as indicated below: \hfill \faFileCodeO \: source code \hfill \faFilePdfO \: PDF file \hfill \faFilePictureO \: image file \hfill \faFileMovieO \: video file

\paragraph*{Question 1 $\vert$ 3 marks}% $\vert$ \faFileCodeO{} \: \faFilePdfO}

Today we're solving the linear advection equation
\begin{subequations}\label{eq:advection}
\begin{align}
\frac{\partial u(t)}{\partial t} + \operatorname{div} (u(t)a(t)) &= 0 && \text{in } Q = ]0,T[ \times \Omega\\
u(0) &= u_0 && \text{in } \Omega\\
u(t) &= g(t) && \text{on } \Sigma_- = \Set{(t,x) \in \: ]0,T[ \times \partial\Omega | a(t,x)\cdot n(x) < 0},
\end{align}
\end{subequations}
where $a: \: ]0,T[ \times \Omega \to \mathbb{R}^2$ is a given vector field, thought of as the flow velocity of a carrier fluid, and $u_0: \Omega \to [0,1]$ the initial concentration of a solute. $g: \: ]0,T[ \times \partial\Omega \to [0,1]$ prescribes the concentration on that part of the boundary $\partial \Omega$ where the flow of solvent is directed into the domain and we use $n$ to denote the outward pointing unit normal vectors on $\partial \Omega$. For simplicity we assume in what follows that $a$ and $g$ do not depend on time.

Recall that upwind discontinuous \textsc{Galerkin} methods employ the bilinear form
\begin{equation*}
\sum_{T\in\mathcal{T}^h} -\int\limits_T ua \cdot \nabla v \dx + \sum_{e\in\mathcal{E}^h} \int\limits_e [v] u_{\text{up}} a \cdot n_+ \ds
\end{equation*}
to discretise the transport term $\operatorname{div} (ua)$.
\begin{enumerate}[(a)]
\item We define the positive and negative parts of a function $f$ by
\begin{equation*}
f^{\text{pos}} = \frac{f + \lvert f \rvert}{2} \qquad \text{and} \qquad f^{\text{neg}} = \frac{f - \lvert f \rvert}{2}.
\end{equation*}
Re-write the edge integral in the upwind DG form in terms of $g$, $u_+$, $u_-$, $(a\cdot n_+)^{\text{pos}}$, $(a\cdot n_+)^{\text{neg}}$ and $[v]$:

\vspace{0.5cm}
\begin{solution}
There are two cases:
\begin{itemize}
\item If $e$ is an interior edge:
\begin{align*}
\int\limits_e [v] u_{\text{up}} a \cdot n_+ \ds
&= \int\limits_e [v] u_{\text{up}} \left[ (a\cdot n_+)^{\text{pos}} + (a\cdot n_+)^{\text{neg}} \right] \ds \\
&= \int\limits_e [v] u_+ (a\cdot n_+)^{\text{pos}} \ds + \int\limits_e [v] u_- (a\cdot n_+)^{\text{neg}} \ds
\end{align*}
The second line follows from the fact that $(a\cdot n_+)^{\text{pos}} = 0$ when $a\cdot n_+ < 0$, and therefore $u_{\text{up}} = u_+$ wherever the integrand in non-zero. Similarly, $u_{\text{up}} = u_-$ in the second term when $(a\cdot n_+)^{\text{neg}}$ is non-zero.
\vfill
\item If $e$ is an exterior edge (no $\pm$ subscripts or $[\:]$ brackets needed, there is only one neighbouring triangle):
\begin{align*}
\int\limits_e [v] u_{\text{up}} a \cdot n_+ \ds
&= \int\limits_e v u_{\text{up}} \left[ (a\cdot n)^{\text{pos}} + (a\cdot n)^{\text{neg}} \right] \ds \\
&= \int\limits_e v u (a\cdot n)^{\text{pos}} \ds + \int\limits_e v g (a\cdot n)^{\text{neg}} \ds
\end{align*}
Where we have used the fact that $u=g$ when $a\cdot n < 0$.
\end{itemize}
\end{solution}

\newpage

\item Complete the \textsf{FEniCS} script \texttt{hw12.py} to solve the linear advection equation with
\begin{equation*}
T = 2\pi \qquad \qquad a(x) = \left(
\begin{array}{c}
-x_2\\ x_1
\end{array}
\right) \textcolor{gray}{ = r \left(
\begin{array}{c}
-\sin \phi \\ \cos \phi
\end{array}
\right)  \text{ in polar coordinates}}  \qquad \qquad  g(x) \equiv 0
\end{equation*}
and the domain and initial data provided. Note that with these data, the advection equation rotates the initial field $u_0$ around the origin in anticlockwise direction, like a rigid body on a turntable.

Use a moderate degree (e.g. $r=1$ or $2$) for the DG-discretisation in space and use the method we are most familiar with, the $\theta$-method, for time stepping. Any value of $\theta$ which results in a stable scheme is fine.

\emph{Hint:} In \textsf{FEniCS} we use \verb|dS| for integrals over interior edges and \verb|ds| for integrals over exterior edges. You must add \verb|('+')| or \verb|('-')| to all discontinuous functions in integral expressions containing \verb|dS|, e.g.
\begin{verbatim}
u('-')*jump(v)*dS
\end{verbatim}
to specify on what side of the edge they should be evaluated. \verb|jump(v)| is a shortcut for \verb|v('+')-v('-')|.

\vspace{0.5cm}
\begin{solution}
See the attached \texttt{hw12.py} code in Appendix~\ref{hw12code}.

Note that \texttt{hw12.py} contains code for both the advection only problem and the advection-diffusion problem:
\begin{itemize}
\item Setting \texttt{D = 0.0} in the script ignores the diffusive flux term and uses the boundary data \texttt{g} on the inflow boundary region $\Gamma_-$
\item Settings \texttt{D = 0.01} (or any non-zero number) ignores the boundary data \texttt{g} and applies no-flux \textsc{Robin} boundary conditions on $\partial\Omega$
\end{itemize}
\end{solution}
\vspace{1cm}

\item Inspect your numerical solution in \textsf{ParaView}, using the `Warp by Scalar'-Filter. Recall that \textsf{FEniCS} generally exports solution data as continuous, piecewise linear functions, so don't be surprised when you cannot see any discontinuous jumps in your visualisation.

Describe any discrepancies between the numerical solution and the exact solution, using the appropriate terminology.

\emph{Hint: } For general triangular meshes, the upwind DG($r$) method discretises the advection term with an accuracy of order $r+\frac{1}{2}$, which increases to $r+1$ on meshes with certain regularity.

\vspace{0.5cm}
\begin{solution}
We know that upwind schemes introduce numerical diffusion into the discretisation of the advection terms.
Thus, we would expect to see numerical diffusion in regions of large concentration gradient, such as at solution discontinuities.
Additionally, although the upwind $DG(r)$ methods preserve mass exactly (to numerical precision), discontinuous changes in the solution can also introduce spurious oscillations, i.e. numerical dissipation.

Indeed, the discontinuous concentration at time $t=0$ begins to spread out as it advects around the domain, and spurious oscillations in the solution occur in the wake behind the advecting concentration mass.
These spurious oscillations become larger with increasing solution degree $r$ and with decreasing amount of implicitness $\theta$ in the time stepping scheme.

This is in contrast to the true (weak) solution with no inflow $g$ which would retain the same discontinuous initial condition and simply advect around the domain in a circular pattern.

\end{solution}

\end{enumerate}

\newpage

\paragraph*{Question 2 $\vert$ 2 marks}% $\vert$ \faFileCodeO{} \:}

Make a copy of your script \texttt{hw12.py} and modify it to solve the advection-diffusion problem
\begin{subequations}\label{eq:advection-diffusion}
\begin{align}
\frac{\partial u(t)}{\partial t} + \operatorname{div} (u(t)a - D\nabla u(t)) &= 0 && \text{in } Q = ]0,T[ \times \Omega\\
u(0) &= u_0 && \text{in } \Omega\\
\left(u(t)a-D\nabla u(t)\right) \cdot n &= 0 && \text{on } \Sigma = ]0,T[ \times \partial\Omega
\end{align}
\end{subequations}
instead, using $D = 0.01$ and all other parameters as in Question 1. The new boundary condition is a no-flux \textsc{Robin} condition. It admits an interpretation of a semi-permeable membrane which allows the solvent, but not the solute to pass through.

Discretise the diffusive flux with the symmetric interior penalty form
\begin{align*}
&\sum_{T\in\mathcal{T}^h} \int\limits_T D\nabla u \cdot \nabla v \dx && \text{(bilinear form of conforming methods)}\\
&- \sum_{e\in\mathcal{E}^h\setminus\mathcal{I}^h} \int_e vD\nabla u \cdot n \ds - \sum_{e\in\mathcal{I}^h} \int_e [v]\left\langle D\nabla u\right\rangle\cdot n_+ \ds && \text{(consistency)}\\
&- \sum_{e\in\mathcal{I}^h} \int_e [u]\left\langle D\nabla v\right\rangle\cdot n_+ \ds && \text{(symmetry)}\\
&+ \sum_{e\in\mathcal{I}^h} \frac{\sigma}{h_e} \int_e [u] [v] \ds && \text{(interior penalty)}
\end{align*}
as derived in class, here written in a form that can be translated directly into \textsf{FEniCS} code (but note that the boundary condition has not been applied yet). Use the penalty parameter $\sigma = 0.1$ and $h_e = \langle h \rangle$.

\emph{Hint:} The commands \verb|h = CellSize(mesh)| and \verb|avg(u)| may be helpful.

\vspace{0.5cm}
\begin{solution}
See the attached \texttt{hw12.py} code in Appendix~\ref{hw12code}.

Note that \texttt{hw12.py} contains code for both the advection only problem and the advection-diffusion problem:
\begin{itemize}
\item Setting \texttt{D = 0.0} in the script ignores the diffusive flux term and uses the boundary data \texttt{g} on the inflow boundary region $\Gamma_-$
\item Settings \texttt{D = 0.01} (or any non-zero number) ignores the boundary data \texttt{g} and applies no-flux \textsc{Robin} boundary conditions on $\partial\Omega$
\end{itemize}
\end{solution}
\vspace{1cm}

\newpage
\paragraph*{Your Learning Progress}% $\vert$ \faFilePdfO}

What is the one most important thing that you have learnt from this assignment?

\vspace{0.5cm}
\begin{itemize}
\item Discontinuous Galerkin methods are really fun to play with and quite easy to implement in \texttt{Fenics}!
\end{itemize}

\vspace{0.8cm}
Any new discoveries or achievements towards the objectives of your course project?

\vspace{0.5cm}
\begin{itemize}
\item Yes, progress has been very good in the last week, especially with having to prepare for the presentation.
I might have to play with discontinuous Galerkin methods for the Bloch-Torrey equation, now!
\end{itemize}

\vspace{0.8cm}
What is the most substantial new insight that you have gained from this course this week? Any \emph{aha moment}?

\vspace{0.5cm}
\begin{itemize}
\item Upwind discretisation schemes in general was kind of one big \emph{aha moment}; I never really understood mathematically why they were used, despite the physical intuition being quite clear.
\end{itemize}

\newpage
\begin{appendices}
% ------------------------------------------------------------ %
% Appendix A: Code
% ------------------------------------------------------------ %
\newpage
\lhead{}
\section{}\label{hw12code}
\textbf{\Large\texttt{hw12.py}}\\\\
\textbf{Note:} unicode characters do not seem to be rendering correctly in \LaTeX; see attached file instead.
\includecode{../fenics/hw12.py}
\end{appendices}

\end{document}
